本文研究机器狗在半径为 1800 米的圆形区域内自动定位并清除无线电干扰源的问题。题目只给出测向误差的有界区间 $[-1^\circ,1^\circ]$，且明确指出同一地点的误差在重复测量下固定不变，因此全文不引入任何概率分布假设，一律以"包含真实目标的可达集合"和"在该集合上取最坏值"的方式建立保证型模型。以此为主线，本文建立了一个由四层递进模型组成的体系：有界测角交会定位、保证接收的第二测点选择、多源滚动式搜索定位清除的联合规划，以及面向未知朝向的覆盖发现证书。

针对问题一，本文将一次示向度观测表示为两个半平面的交集，即半宽 $\varepsilon=1^\circ$ 的正向角域，从而把定位区域写成闭凸多面集 $P=\{x:Ax\le b\}$。该表示不使用斜率，在 $90^\circ$、$270^\circ$ 及 $0^\circ/360^\circ$ 跨界附近仍然一致。区域状态（空集、无界、点、线段、二维多边形）通过四个方向上的线性规划判定，无需人为设置裁剪框；有界时枚举不平行边界直线对的交点并逐点回代，得到顶点集。区域直径由顶点对取得，即 $D(P)=\max_{i,j}\lVert v_i-v_j\rVert$，其证明基于凸组合与三角不等式。针对"以定位区域直径为直径的圆能否覆盖该区域"这一问题，本文证明该圆存在当且仅当最小包围圆半径 $\varrho(P)=D(P)/2$，并进一步给出一个可由本题测向机制真实实现的等边三角形反例：三条下边界半平面恰好围成边长 38 米的三角形，其 $\varrho=38/\sqrt{3}\approx21.94$ 米，故直径 40 米的圆同样无法覆盖。由此得到与后续两问衔接的结论：单个位置 $q$ 保证清除的充要条件是 $\max_i\lVert v_i-q\rVert\le20$，所有保证清除位置构成等半径圆盘的交 $\mathcal{C}(P)=\bigcap_i\overline{B}(v_i,20)$。

针对问题二，本文首先指出：首次返回有效示向度这一事实本身蕴含真实距离 $r\le R$（$R$ 为该源未知的接收半径，$1000\le R\le1500$ 米）。利用这一关系，本文推导出无需知道目标距离与朝向、即可保证第二检测点 $q$ 仍能收到信号的充分区域——三个半径均为 $R_0=1000$ 米的圆盘之交，其局部坐标形式为 $a^2+b^2\le R_0^2$ 且 $a^2+b^2\le2R_0(a\cos\varepsilon-\lvert b\rvert\sin\varepsilon)$。该结论同时说明单纯横向平移并不可靠：当真实目标位于 $(1000,0)$ 时横移至 $(0,500)$ 将失去信号。此后以横向分离阈值排除与目标共线的退化观测，并以第二次观测后可行域的最坏最小包围圆半径 $F(q)=\sup_{\psi}\varrho(P_2(q,\psi))$ 作为评价指标，在移动预算下给出策略族。对首次位于原点、读数 $0^\circ$ 的示例，移动预算 250、500、750、1000 米对应的推荐点最坏后验半径上界依次为 220.01、123.93、81.44、56.21 米。需要说明的是，这是有限候选搜索的结果，而非连续空间全局最优；且在 1000 米预算下仍存在超过 20 米的实际合成后验，故本文不宣称增加一次观测即可普遍保证清除。

针对问题三，本文把整场任务的虚拟时间写成可加账本 $T=L/5+5M+S+5C+3F$，并在此基础上建立"发现—定位—清除—收尾"闭环。定位环节直接复用问题一的半平面交与最小包围圆；负观测则以"距离超过接收半径"与"同一源的接收半径恒定"两个事实，转化为排除圆盘与正负观测对生成的半平面。搜索环节的关键是给出可检验的终止依据：把场地划为 $128\times128$ 方格，只有当代方格全部角点到某次实际无信号位置的距离都不超过 999.99 米时才排除该格，全部保留格被覆盖即证明该频道不存在；此外允许清除到题目上限 16 个源后停止。路线采用固定起点、自由终点的最近邻加 2-opt 滚动重规划。在 200 场同场配对审计中，全部场景均完整清除、零失败清除，场景等权平均由 317.38 秒/点降至 237.19 秒/点（降低 25.27\%）。在随后一次真实官方演练（案例 JCEK-BMVS-WD8Z-WF6J）中，官方确认场内 14 个全向源，本策略清除 14 个、零失败，总虚拟时间 3854.08 秒、平均 275.29 秒/源，累计移动 16020.42 米，其中移动占 83.13\%。

针对问题四，本文首先论证第三问的两项核心机制在此失效：其一，无信号不再能排除 1000 米圆盘，因为定向源可能就在近处背向接收机；其二，第三问的全向覆盖终止证书不再成立。为此本文给出一项与朝向无关的发现证书：用边长 990 米的等边三角网覆盖场地，若某频道在某个三角形的三个顶点都实际测得无信号，则由凸组合 $\sum_i\lambda_in^{T}(s_i-g)=0$ 可知任意朝向下至少有一个顶点位于源的闭发射半平面内，且该顶点到源的距离不超过 990 米，小于接收半径下界 1000 米，故该三角形内必无此频道的源。在此基础上，本文进一步把覆盖站点的构造一般化为凸包排除条件：对每个凸子区域 $Q$，若存在一组到 $Q$ 内所有点距离都不超过 1000 米的实际无信号站点，且 $Q$ 包含于这些站点的凸包，则 $Q$ 中不可能存在任意朝向的定向源。据此构造出中心 1 点、半径 980 米内圈 7 点、半径 1850 米外圈 14 点共 22 个站点的布局，替代原先的 31 点网格，并用 1334 个凸子区域给出全场覆盖证据。40 场同场配对中全部完整清除，平均由 696.78 秒/点降至 515.63 秒/点（降低 26.00\%），39 场改善、1 场退化。

最后，本文用独立几何核验、消融实验与压力测试检验上述模型。问题一的独立验证包含 260 项检查，用线性规划支持函数核对顶点枚举（最大差值约 $4.55\times10^{-13}$ 米），用缩放凸范数优化核对最小包围圆（半径差约 $5.96\times10^{-10}$ 米）；问题二包含 1305 项检查与 1100 条合成方向工况。需要特别指出的是，第三、四问的全部数值均为本地合成场景结果，尚未完成题目要求的三次正式测试，相应结果表在本文中留空；第三问的移动下界分析（圆周 10 源）给出 233.64 秒/点，说明 200 秒/点这一内部研究目标无法在所有合法场景上获得硬保证，且当前既未达到该目标，也未能证明其分布均值不可达。
